
% --------------------------------------------------------------------------------------------------
% Kommando Beispiele
% --------------------------------------------------------------------------------------------------

% Referenz auf z.B. Kapitel
\chapter{Comparison of Proof Assistants}
\label{c:chapter6}

Here the next chapter \ref{c:chapter6}...

\subsection{Is Lean Just Like Rocq?}

Ref: lean FAG rcoc. 
On a foundational level, Lean and \href{https://rocq-prover.org/}{Rocq} 
(the theorem prover formerly known as Coq) have much in common: both 
are based on dependent type theory (the Calculus of Inductive 
Constructions), with some technical differences:

\begin{itemize}
    \item Lean uses definitional proof irrelevance.
    \item Lean uses a non-cumulative hierarchy of universes, while 
    Rocq uses a cumulative hierarchy.
    \item Rocq's basic theory supports recursion, and its kernel has 
    a trusted termination checker, whereas Lean's basic theory only 
    provides recursors, and the handling of recursion is not part 
    of the trusted code base.
\end{itemize}

On top of these similar foundations, however, the systems are 
engineered rather differently. Lean's most important design goals 
are scalability, automation, and a focus on classical logic. Lean 
is also known for its excellent metaprogramming framework, in which 
Lean itself is used for metaprogramming and extensibility. Lean is 
also written in Lean, and many users have written their own 
extensions. Furthermore, Lean provides a rich and modern user 
interface.